% [原文 p.71]
% 块首承接：p.70 末行起的跨页句（“……这说明：单射同态……”）已由 ch02d 收尾，本块自下一段起。
到此为止所叙述的内容，对有限群 $G$ 的一切 $G$ 空间同样成立，本节到此结束。此处用到的只是同态延拓一类基本事实，它们在 $G$ 向量丛的情形照样成立。因此，$K_G(X,Y)$ 的元素可以用 $X$ 上向量丛的 $G$ 复形来表示，而该复形在 $Y$ 上非轮状。特别地，$G$ 向量丛 $V$ 的外积代数确定元素
\[
\lambda_V\in\tilde K_G(X^V).
\]

\section{Thom 同构}
\takeaway{本节把 §2.5 的射影丛计算升级为 Thom 同型 $K^*(X)\cong\tilde K^*(X^E)$，实现手段是乘以 Thom 类 $\lambda_E$；路线是分裂原理：先在可分解的 $E$ 上算，再用旗丛的单射拉回一般情形。}

若 $E=\sum L_i$ 是 $X$ 上的可分解向量丛（即线丛的直和），则由命题 2.5.3 可以把 $K(P(E))$ 的结构作为 $K(X)$ 代数确定下来。现在，对任意空间 $X$ 有自然同构映射
\[
K^*(X)\cong K(X\times S^1).
\]
另外，若 $\pi\colon X\times S^1\to X$ 是投影，则
\[
P(E)\times S^1=P(\pi^*E),
\]
因此
\[
K^*(P(E))\cong K(P(\pi^*E)).
\]
于是，把命题 2.5.3 中的 $X$ 换成 $X\times S^1$，立即得到下面的命题。

% [原文 p.72]
\begin{proposition*}{2.7.1}
  设 $E=\sum L_i$ 为 $X$ 上的可分解向量丛。此时，作为 $K^*(X)$ 代数，$K^*(P(E))$ 由 $[H]$ 生成，并服从唯一的关系式
  \[
  \prod_i\bigl([L_i][H]-1\bigr)=0.
  \]
\end{proposition*}

\begin{remark}{}
  这一点与命题 2.5.3 一样，可立即推广到 $G$ 空间，给出 $K_G^*(P(E))$ 作为 $K_G^*(X)$ 代数的表示。
\end{remark}

现在，若把 Thom 空间 $X^E$ 与 $P(E\oplus1)/P(E)$ 等同，则由 2.6 节末尾可知，$\lambda_E$ 在 $K(P(E\oplus1))$ 中的像是
\[
\sum(-1)^i[H]^i[\lambda^iE]=\prod_i\bigl(1-[L_i][H]\bigr).
\]
这个元素生成
\[
K^*(P(E\oplus1))\longrightarrow K^*(P(E))
\]
的核（作为理想），于是得到下面的命题。

\begin{proposition*}{2.7.2}
  设 $E$ 为 $X$ 上的可分解向量丛。此时，$\tilde K^*(X^E)$ 是以 $\lambda_E$ 为生成元的自由 $K^*(X)$ 模。
\end{proposition*}

\begin{remark}{}
  这种可分解情形的“Thom 同型定理”，与前面一样对 $G$ 空间也成立。下面说明这个事实可以怎样使用。
\end{remark}

\begin{corollary*}{2.7.3}
  设 $X$ 为 $G$ 空间，且 $K_G^1(X)=0$。设 $E$ 为可分解的 $G$ 向量丛。此时，若 $S(E)$ 是球面丛，则有正合列
  \[
  0\longrightarrow K_G^1(S(E))\longrightarrow K_G^0(X)\xrightarrow{\ \varphi\ }K_G^0(X)
  \longrightarrow K_G^0(S(E))\longrightarrow0,
  \]
  其中 $\varphi$ 是乘以
  \[
  \lambda_{-1}[E]=\sum(-1)^i\lambda^i[E]
  \]
  的映射。
\end{corollary*}

% [原文 p.73]
\begin{proof}
  把命题 2.7.2 用于对 $(B(E),S(E))$ 的正合列，便可立即得到。
\end{proof}

% [Tutor Note: 点拨：命题 2.7.1 的关系式 $\prod_i([L_i][H]-1)=0$ 与命题 2.5.3 的 $\prod_i(t-[L_i^*])$ 是同一个关系：令 $t=[H]^{-1}=[H^*]$ 再乘以 $[H]^n$ 即可互推。抄录时注意两式中 $1$ 与 $[H]$ 的位置不同。]

为把这个推论用于 $X$ 是单点的情形，下面给出一个引理。

\begin{lemma*}{2.7.4}
  $K_G^1(*)=0$。
\end{lemma*}

\begin{proof}
  只需证明映射
  \[
  K_G(S^1)\longrightarrow K_G(*)
  \]
  是同构即可。然而 $G$ 平凡地作用在 $S^1$ 上，所以
  \[
  \begin{aligned}
  K_G(S^1)&\cong K(S^1)\otimes R(G)\\
  &\cong K(*)\otimes R(G)\\
  &\cong K_G(*)
  \end{aligned}
  \]
  成立。
\end{proof}

于是，可以在推论 2.7.3 中取 $X$ 为单点。进一步，取 $G$ 为阿贝尔群，则 $E$ 总是可分解的。因此得到下面的推论。

\begin{corollary*}{2.7.5}
  设 $G$ 为阿贝尔群，$E$ 为 $G$ 模。此时有正合列
  \[
  0\longrightarrow K_G^1(S(E))\longrightarrow R(G)\xrightarrow{\ \varphi\ }R(G)
  \longrightarrow K_G^0(S(E))\longrightarrow0,
  \]
  其中 $\varphi$ 是乘以
  \[
  \lambda_{-1}[E]=\sum(-1)^i\lambda^i[E]
  \]
  的映射。
\end{corollary*}

特别地，若 $G$ 自由地作用在 $S(E)$ 上（此时 $G$ 必然是循环群），则
% [原文 p.74]
\[
K_G^*(S(E))\cong K^*(S(E)/G)
\]
成立。这样，便得到下面的推论。

\begin{corollary*}{2.7.6}
  设 $G$ 为循环群，$E$ 为 $G$ 模，且 $G$ 对 $S(E)$ 的作用是自由的。此时有正合列
  \[
  0\longrightarrow K^1(S(E)/G)\longrightarrow R(G)\xrightarrow{\ \varphi\ }R(G)
  \longrightarrow K^0(S(E)/G)\longrightarrow0,
  \]
  其中 $\varphi$ 是乘以 $\lambda_{-1}[E]$ 的映射。
\end{corollary*}

\begin{remark}{}
  当有限循环群以外的群 $G$ 自由地作用在球面上时，只要能把 $K_G$ 的 Thom 同型推广到未必可分解的向量丛，同样的结果也成立。不过，这不在本讲义中讨论。
\end{remark}

作为推论 2.7.6 的特殊情形，考虑 $G=\Z/2$ 非平凡地作用在 $E=\C^n$ 上。此时
\[
S(E)/G=P_{2n-1}(\R)
\]
是奇数次维数的实射影空间。
\[
R(\Z/2)=\Z[\rho]/(\rho^2-1),\qquad \lambda_{-1}[E]=(1-\rho)^n,
\]
而令 $\sigma=\rho-1$，则 $\sigma^2=-2\sigma$ 且 $\lambda_{-1}[E]=(-\sigma)^n$ 成立。所以可知 $\tilde K^0(P_{2n-1}(\R))$ 是阶为 $2^{n-1}$ 的循环群，而 $K^1(P_{2n-1}(\R))$ 是无限循环群。比较 $n$ 与 $n+1$ 对应的正合列，得到交换图
\[
\begin{array}{ccccccc}
0 & \longrightarrow & K^1(P_{2n+1}(\R)) & \longrightarrow & R(\Z/2)
& \xrightarrow{(-\sigma)^{n+1}} & R(\Z/2)\\[2pt]
& & \big\downarrow & & \big\downarrow{\scriptstyle-\sigma} & & \big\downarrow{\scriptstyle 1}\\[2pt]
0 & \longrightarrow & K^1(P_{2n-1}(\R)) & \longrightarrow & R(\Z/2)
& \xrightarrow{(-\sigma)^{n}} & R(\Z/2)
\end{array}
\]
% [原文 p.75]
然而，在 $R(\Z/2)$ 中，对 $n\ge1$，$(-\sigma)^n$ 的核与 $-\sigma$ 的核一致。因此，映射
\[
K^1(P_{2n+1}(\R))\longrightarrow K^1(P_{2n-1}(\R))
\]
是零映射。由对 $(P_{2n+1}(\R),P_{2n}(\R))$ 与 $(P_{2n}(\R),P_{2n-1}(\R))$ 的正合列可知，
\[
K^1(P_{2n+1}(\R))\longrightarrow K^1(P_{2n}(\R))
\]
是满射，并且
\[
K^1(P_{2n}(\R))\longrightarrow K^1(P_{2n-1}(\R))
\]
是单射。因此
\[
K^1(P_{2n}(\R))=0.
\]
于是，由对 $(P_{2n+1}(\R),P_{2n}(\R))$ 的正合列可知，
\[
K^0(P_{2n+1}(\R))\longrightarrow K^0(P_{2n}(\R))
\]
是同构。总结起来，下面的命题得到证明。

% [Tutor Note: 易错预警：$R(\Z/2)=\Z[\rho]/(\rho^2-1)$ 中 $\rho$ 是自然表示的特征标，$\sigma=\rho-1$ 对应一维非平凡表示；关系 $\sigma^2=-2\sigma$ 表明 $\Z[\rho]/(\rho^2-1)$ 的两个生成元之和与差都给出 $\Z$ 的直和项。$\lambda_{-1}[E]=(-\sigma)^n$ 的形式决定下文 $2^{n-1}$ 的阶，符号不能丢。]

\begin{proposition*}{2.7.7}
  $K^*(P_n(\R))$ 的结构如下：
  \[
  K^1(P_{2n+1}(\R))=\Z,\qquad K^1(P_{2n}(\R))=0,
  \]
  \[
  \tilde K^0(P_{2n+1}(\R))=\tilde K^0(P_{2n}(\R))=\Z/2^n.
  \]
\end{proposition*}

把推论 2.7.6 用于其他空间，作为练习题留给读者。

现在，转向一般的 Thom 同型定理。这里要强调，所用的手法不能推广到 $G$ 向量丛。推广到 $G$ 向量丛需要完全不同的手法，此处不讨论。

从下面的一般性结果开始。

\begin{theorem*}{2.7.8}
  设 $\pi\colon B\to X$ 为紧空间之间的映射，$\mu_1,\ldots,\mu_n$ 为 $K^*(B)$ 的齐次元素。设 $M^*$ 为由 $\mu_1,\ldots,\mu_n$ 生成的自由 $(\Z/2)$ 分次群。若对每一点 $x\in X$ 都存在邻域 $U$，使得对一切 $V\subset U$，自然映射
  \[
  K^*(V)\otimes M^*\longrightarrow K^*(\pi^{-1}(V))
  \]
  都是同构映射，则对任意 $Y\subset X$，映射
  \[
  K^*(X,Y)\otimes M^*\longrightarrow K^*(B,\pi^{-1}(Y))
  \]
  是同构映射。
\end{theorem*}

% [原文 p.76]
\begin{proof}
  若 $U\subset X$ 具有“对一切 $V\subset U$ 都有
  \begin{equation}
  K^*(V)\otimes M^*\cong K^*(\pi^{-1}(V)) \tag{1}
  \end{equation}
  成立”这一性质，就说 $U$ 是良好（good）的。若 $U$ 良好，利用正合列与“$\otimes M^*$ 保持正合性（$M^*$ 无挠）”这一事实，可以导出
  \begin{equation}
  K^*(U,V)\otimes M^*\cong K^*(\pi^{-1}(U),\pi^{-1}(V)). \tag{2}
  \end{equation}
  这里当然用到了 $\delta$ 与乘法相容这一点（引理 2.6.1）。因此，必须证明的是
  \[
  X\ \text{局部良好}\ \Longrightarrow\ X\ \text{良好}.
  \]
  由于 $X$ 紧，只需证明
  \[
  U_1,U_2\ \text{良好}\ \Longrightarrow\ U_1\cup U_2\ \text{良好}.
  \]
  现在，任意 $V\subset U_1\cup U_2$ 都形如 $V=V_1\cup V_2$，其中 $V_i\subset U_i$（从而 $V_i$ 也良好）。由于
  \[
  V/V_2=V_1/(V_1\cap V_2),
  \]
  可知对 $(V,V_2)$ 式 (2) 成立。对 $V_2$ 式 (1) 成立，所以由 $(V,V_2)$ 的正合列，式 (1) 对 $V$ 成立。因此，
% [原文 p.77]
  $U_1\cup U_2$ 良好，证明完成。
\end{proof}

% [Tutor Note: 点拨：定理 2.7.8 的论证骨架是“局部同构 $\Rightarrow$ 两开集之并仍良好 $\Rightarrow$ 由紧性得整体”。中间那一步用对 $(V,V_2)$ 的正合列（Mayer–Vietoris）把 $K^*(V)$ 夹在中间；$M^*$ 无挠保证 $\otimes M^*$ 保持单射，此条件不可省。]

\begin{corollary*}{2.7.9}
  设 $\pi\colon E\to X$ 为 $n$ 维向量丛，$H$ 为 $P(E)$ 上的万有丛。此时，$K^*(P(E))$ 是以 $1,[H],[H]^2,\ldots,[H]^{n-1}$ 为基的自由 $K^*(X)$ 模。$[H]$ 满足关系式
  \[
  \sum(-1)^i[H]^i[\lambda^iE]=0.
  \]
\end{corollary*}

\begin{proof}
  $E$ 是局部平凡的，因而特别地是局部可分解的\footnote{这一点不能推广到 $G$ 空间。}。于是，由命题 2.7.1，对每一点 $x\in X$ 存在邻域 $U$，使得对一切 $V\subset U$，$K^*(P(E|_V))$ 是以 $1,[H],\ldots,[H]^{n-1}$ 为基的自由 $K^*(V)$ 模。这里应用定理 2.7.8。$[H]$ 应满足的关系式已在 2.6 节末尾证明。
\end{proof}

\begin{corollary*}{2.7.10}
  设 $\pi\colon E\to X$ 为向量丛。以 $F(E)$ 表示 $E$ 的旗丛，其投影为 $p\colon F(E)\to X$。此时，$p^*\colon K^*(X)\to K^*(F(E))$ 是单射。
\end{corollary*}

\begin{proof}
  $F(E)$ 是射影丛 $P(E)$ 上的旗丛，而后者的维数比 $\dim E$ 小 1。用推论 2.7.9 对 $\dim E$ 作归纳即可证明。
\end{proof}

\begin{corollary*}{2.7.11}（分裂原理）
  对 $X$ 上的向量丛 $E_1,\ldots,E_n$，存在满足下列条件的空间 $F$ 与映射 $\pi\colon F\to X$。
  \begin{enumerate}[label=\arabic*),leftmargin=2.2em,itemsep=2pt]
  \item $\pi^*\colon K^*(X)\to K^*(F)$ 是单射。
  \item 各个 $\pi^*(E_i)$ 都是线丛的直和。
  \end{enumerate}
\end{corollary*}

\begin{proof}
  取 $F$ 为 $\bigoplus_i E_i$ 的旗丛即可。
\end{proof}

分裂原理的重要性是显然的。借助分裂原理，许多问题都可以归结为可分解向量丛的情形。

\begin{corollary*}{2.7.12}（Thom 同型定理）
  若 $\pi\colon E\to X$ 是向量丛，则由 $\Phi(x)=\lambda_E x$ 定义的
% [原文 p.78]
  \[
  \Phi\colon K^*(X)\longrightarrow\tilde K^*(X^E)
  \]
  是同构映射。
\end{corollary*}

\begin{proof}
  与命题 2.7.2 由命题 2.7.1 导出的方式相同，这个推论由推论 2.7.9 导出。
\end{proof}

% [Tutor Note: 点拨：分裂原理之所以够用，是因为 $\pi^*$ 是单射：一个等式在 $K^*(X)$ 中成立当且仅当它在 $K^*(F)$ 中成立，而 $\pi^*(E_i)$ 已分解为线丛之和，命题 2.7.1、2.7.2 便可逐项使用。这也解释了作者为何先算可分解情形。]

下面的命题 2.7.13 与 2.7.14 的证明，作为练习题留给读者。

\begin{proposition*}{2.7.13}
  设 $\pi\colon E\to X$ 为向量丛，且 $L_1,\ldots,L_n$ 是 $F(E)$ 上的万有线丛。此时，由 $t_i\mapsto[L_i]$ 定义的映射是 $K^*(X)$ 模的同构映射
  \[
  K^*(X)[t_1,\ldots,t_n]/I\longrightarrow K^*(F(E)),
  \]
  其中 $I$ 是如下元素生成的理想：以 $\sigma^i$ 记 $i$ 次基本对称式，则生成元为
  \[
  \sigma^1(t_1,\ldots,t_n)-E,\quad
  \sigma^2(t_1,\ldots,t_n)-\lambda^2(E),\quad\ldots,\quad
  \sigma^n(t_1,\ldots,t_n)-\lambda^n(E).
  \]
\end{proposition*}

\begin{proposition*}{2.7.14}
  设 $\pi\colon E\to X$ 为 $n$ 维向量丛，$G_k(E)\xrightarrow{\ p\ }X$ 为 $E$ 的（$k$ 维子空间的）Grassmann 丛。设 $F$ 为 $G_k(E)$ 上诱导出的 $k$ 维万有丛，$F'$ 为商丛 $p^*(E)/F$。此时，$t_i\mapsto\lambda^i(F)$、$s_i\mapsto\lambda^i(F')$ 是 $K^*(X)$ 模的同构映射
  \[
  K^*(X)[t_1,\ldots,t_k,s_1,\ldots,s_{n-k}]/I\longrightarrow K^*(G_k(E)),
  \]
  其中 $I$ 是由元素
  \[
  \Bigl(\sum_{i+j=\ell}t_is_j\Bigr)-\lambda^\ell(E)\qquad(\ell\ge0)
  \]
  生成的理想（提示：把 $G_k(E)$ 与 $E$ 的旗丛作比较）。
\end{proposition*}

特别地，若 $G_{n,k}$ 是 $n$ 维向量空间的 $k$ 维子空间构成的 Grassmann 流形，则 $K^*(G_{n,k})$ 无挠。这也可以由 $G_{n,k}$ 的胞腔分解导出。由归纳法可知，对 Grassmann 流形的乘积空间，其 $K^*$ 也无挠。

% [原文 p.79]
\begin{theorem*}{2.7.15}
  设 $X$ 为使 $K^*(X)$ 无挠的空间，$Y$ 为（有限）胞腔复形，$Y'\subset Y$ 为子复形。此时，映射
  \[
  K^*(X)\otimes K^*(Y,Y')\longrightarrow K^*(X\times Y,X\times Y')
  \]
  是同构映射。
\end{theorem*}

\begin{proof}
  当 $Y$ 是球体、$Y'$ 是其边界时，这个定理是命题 2.7.2 的推论。因此，对 $Y$ 所含胞腔的个数作归纳，可知定理对任意 $(Y,Y')$ 成立。
\end{proof}

\begin{corollary*}{2.7.16}（Künneth 定理）
  设 $X$ 为使 $K^*(X)$ 是有限生成阿贝尔群的空间，$Y$ 为胞腔复形。此时有自然正合列
  \[
  \begin{aligned}
  0\longrightarrow\sum_{i+j=k}K^i(X)\otimes K^j(X)&\longrightarrow K^k(X\times Y)\\
  &\longrightarrow\sum_{i+j=k+1}\Tor\bigl(K^i(X),K^j(X)\bigr)\longrightarrow0,
  \end{aligned}
  \]
  其中所有下标都取 $\Z/2$ 的元素。
\end{corollary*}

% [Tutor Note: 原书 p.79 把 Künneth 公式中的 $K^j(Y)$ 印成了 $K^j(X)$（两处都是）。按公式的含义应为 $K^j(Y)$：左端是 $K^*(X)\otimes K^*(Y)\to K^*(X\times Y)$ 的 Künneth 分解。正文照原书抄录，请按 $K^j(Y)$ 理解。]

\begin{proof}
  设能找到空间 $Z$ 与嵌入 $f\colon X\to Z$，使 $K^*(Z)$ 无挠且 $f^*\colon K^*(Z)\to K^*(X)$ 是满射。此时由对的正合列可知，$K^*(Z/X)$ 无挠。由定理 2.7.15，$K^*(Z\times Y)=K^*(Z)\otimes K^*(Y)$，$K^*((Z/X)\times Y)=K^*(Z/X)\otimes K^*(Y)$。于是，推论的结果由对 $(Z\times Y,X\times Y)$ 的正合列导出。

  下面构造这样的映射 $g\colon SX\to Z$。设 $a_1,\ldots,a_n$ 生成 $K^0(X)$，$b_1,\ldots,b_m$ 生成 $K^{-1}(X)=K(SX)$。此时，每个 $a_i$ 对适当的 $r_i$、$s_i$ 确定映射 $\alpha_i\colon X\to G_{r_i,s_i}$，每个 $b_i$ 确定映射 $\beta_i\colon SX\to G_{u_i,v_i}$。令 $\alpha\colon X\to G_{r_1,s_1}\times\cdots\times G_{r_n,s_n}=G'$ 为 $\alpha_1\times\cdots\times\alpha_n$，令 $\beta\colon SX\to G_{u_1,v_1}\times\cdots\times G_{u_m,v_m}=G''$ 为 $\beta_1\times\cdots\times\beta_m$，则
% [原文 p.80]
  \[
  \alpha^*\colon K^0(G')\longrightarrow K^0(X)\ \text{是满射},
  \]
  \[
  \beta^*\colon K^0(G'')\longrightarrow K^0(SX)\ \text{是满射}.
  \]
  因此，令 $f\colon(S\alpha)\times\beta\colon SX\to(SG')\times G''=G$，则
  \[
  f^*\colon K^*(G)\longrightarrow K^*(SX)\ \text{是满射},
  \]
  而 $K^*(G)$ 如所求那样无挠。这样，关于 $SX$ 的式子得到证明，而它与关于 $X$ 的式子等价。
\end{proof}

% [Tutor Note: 次数约定：原书此处写 $K^{-1}(X)=K(SX)$；按定义 2.4.1，$K^{-1}(X)=\tilde K(S(X^+))$，故右端应理解为约化群 $\tilde K(SX)$。$K(SX)$ 还多一个 $\Z$ 的直和项，读时按约化群处理。]

下面，用 $U(n)$ 表示 $n$ 阶酉群，计算环 $K^*(U(n))$。对任意紧 Lie 群 $G$，可以认为表示 $\rho\colon G\to\mathrm{GL}(m,\C)$ 定义了元素 $[\rho]\in K^1(G)$。为此，只需把 $\rho$ 看作映射并忽略其乘法性即可。这里，设 $\alpha,\beta\colon G\to\mathrm{GL}(m,\C)$ 是两个表示，且它们在闭子群 $H$ 上一致。此时可以定义映射
\[
\gamma\colon G/H\longrightarrow\mathrm{GL}(m,\C)
\]
为 $\gamma(gH)=\alpha(g)\beta(g)^{-1}$。由于 $\alpha$、$\beta$ 的乘法性，这是无矛盾地定义好的。映射 $\gamma$ 定义了元素 $[\gamma]\in K^1(G/H)$，而它在 $K^1(G)$ 中的像恰好是 $[\alpha]-[\beta]$。作为其特殊情形，考虑
\[
G=U(n),\qquad H=U(n-1),\qquad G/H=S^{2n-1}.
\]
设 $\alpha$、$\beta$ 分别是外积代数 $\Lambda^*(\C^n)$ 的偶数次部分与奇数次部分所给出的 $G$ 的表示，并通过与 $\C^n$ 的第 $n$ 个基向量 $e_n$ 作外积 $\cdot\wedge e_n$ 把这两部分等同起来。$U(n-1)$ 不动 $e_n$，所以这个等同与 $U(n-1)$ 的作用相容。这样便得到元素
\[
[\gamma]\in K^1(S^{2n-1}).
\]
把它移到与之同构的群 $\tilde K(S^{2n})$ 中，则由其定义可知，$[\gamma]$ 正是 2.6 节用外积代数构造的基本元素
% [原文 p.81]
\[
\lambda_{\C^n}\in\tilde K(S^{2n})
\]
本身。因此，$[\gamma]$ 是 $K^1(S^{2n-1})$ 的生成元，而它在 $K^1(U(n))$ 中的像是 $\sum(-1)^i[\lambda^i]$。其中 $\lambda^i$ 是外积幂表示。这些预备考察为证明下面的定理作好了准备。

\begin{theorem*}{2.7.17}
  设 $\lambda^i$ 为 $U(n)$ 的 $i$ 次外积幂表示。此时，$K^*(U(n))$ 是由 $[\lambda^1],\ldots,[\lambda^n]$ 生成的外积代数。
\end{theorem*}

\begin{proof}
  对 $n$ 用归纳法证明。考虑如下映射：
  \[
  U(n)\longrightarrow U(n)/U(n-1)=S^{2n-1}.
  \]
  用 $\mu^i$ 表示 $U(n-1)$ 的 $i$ 次外积幂表示，则 $\lambda^i$ 到 $U(n-1)$ 的限制是 $\mu^i\otimes\mu^{i-1}$，所以把归纳法假设与定理 2.7.8 结合起来，可知 $K^*(U(n))$ 是由 $[\lambda^1],\ldots,[\lambda^{n-1}]$ 构成的单项式生成的自由 $K^*(S^{2n-1})$ 模。然而 $K^*(S^{2n-1})$ 是由一个生成元 $[\gamma]$ 生成的外积代数，而它在 $K^*(U(n))$ 中的像如前所示，为
  \[
  \sum_{i=0}^{n}(-1)^i[\lambda^i].
  \]
  结果，$K^*(U(n))$ 如所求，是 $[\lambda^1],\ldots,[\lambda^n]$ 上的外积代数。
\end{proof}

% [Tutor Note: 点拨：$\lambda^i$ 到 $U(n-1)$ 的限制为 $\mu^i\otimes\mu^{i-1}$，这来自外积代数的分裂 $\Lambda^*(\C^n)=\Lambda^*(\C^{n-1})\oplus e_n\Lambda^*(\C^{n-1})$：含 $e_n$ 的那一半把次数升 1，正是 $\mu^{i-1}$ 乘以 $\mu^1$ 的效果。这一步是归纳法能启动的关键。]

\chapterreview{第 2 章是全书的重心，也是 K 理论真正开始运转的地方。它把第 1 章关于向量丛的分类结果（同伦分类、衔接函数、Grassmann 流形与旗流形的分类定理）组装成一个可计算的上同调论，并立起两根支柱：Bott 周期性定理与 Thom 同型定理。第 1 章只回答“向量丛有多少种”，第 2 章回答“这些类如何运算”，由此得到的 $K^*(X)$ 是分次交换环，乘法来自向量丛的张量积。此后第 3 章的外幂算子、附录 B 的幂算子、附录 C 的实结构与系数理论，都建立在这个环及其运算之上。

全章的逻辑推进可分四步，四步之间靠“把难算的对象搬到好算的对象上”这一策略串起来。第一步（§2.1–§2.2）从 $K(X)$ 的半群构造出发，把周期性定理化归为射影丛 $P(L\oplus1)$ 上的等式 $([H]-[1])([L][H]-[1])=0$，再用 $K(X)[t]/((t-1)([L]t-1))\cong K(P)$ 把整个定理压缩成一个代数同构；证明依靠 Fourier 级数的 Cesàro 平均与线性衔接函数的谱分解。第二步（§2.3–§2.4）把 $K$ 提升为定义在紧空间对上的上同调论，用约化悬垂与 smash 积定义负次数群，得到 6 项正合列、周期同构 $\beta$ 与 Künneth 公式，并随即指出这些证明是形式化的，对任何满足公理的一般上同调论都成立。第三步（§2.5）落实计算：球面、胞腔复形、射影丛、Grassmann 流形与旗流形逐一算出，Brouwer 不动点定理作为副产品出现。第四步（§2.6–§2.7）引入相对群中的乘法与 Thom 类 $\lambda_E$，用分裂原理把一般向量丛拉到线丛的和上，最后以 Thom 同型收束全章。

需要记住的结论有六条：Bott 周期性 $\tilde K(S^2X)\cong\tilde K(X)$；分次环 $K^*(X)=K^0(X)\oplus K^1(X)$ 与 $K^n=K^{n-2}$ 的同一化；$\tilde K(S^n)$ 奇零偶 $\Z$，由此得到 Brouwer 不动点定理；射影丛公式 $K^*(P(E))\cong K^*(X)[t]/\prod_i(t-[L_i^*])$；分裂原理把一般向量丛的问题归约到线丛之和；Thom 同型 $K^*(X)\cong\tilde K^*(X^E)$ 与乘法性 $\lambda_{V\oplus W}=\lambda_V\lambda_W$。最后一条是全章的技术顶点，也是第 3 章计算外幂算子的唯一入口。

几处易混之处值得重申。$K(X,Y)$ 的定义是 $\tilde K(X/Y)$，而不是 $K(X)$ 对 $K(Y)$ 的商，只有 $Y\to X$ 是好对时两者才一致。$K^n$ 与 $K^{n-2}$ 始终同一化，6 项正合列才能首尾相接成环，否则 $\delta$ 的次数无法闭合。$H$ 与其对偶 $H^{-1}$、$L$ 与 $L^*$ 在各公式中反复换写，务必用 $[H][H^{-1}]=[1]$ 统一。$\lambda_E$ 的符号约定（外积代数取交替和 $\sum(-1)^i[H]^i[\lambda^iE]$）直接决定 Thom 同构的符号，也决定了实射影空间 $P_n(\R)$ 中那个 $2^{n-1}$ 的阶。此外，$G$ 等变版本不能照搬：§2.7 开头那句“所用的手法不能推广到 $G$ 向量丛”是全章最重要的限制性说明。

从今天的观点看，第 2 章给出的是 Atiyah–Hirzebruch 谱序列的特殊情形，但作者刻意绕开谱序列，改用射影丛与分裂原理做初等化处理。这条路线对初学者更透明，也解释了全章几乎无图、却处处是交换图式论证的原因：全部信息都压在 $\lambda_E$ 与 $[H]$ 这两个类上。$G$ 等变情形只在 §2.3 与 §2.7 的注中点到为止——可分解向量丛上的 Thom 同型能够推广，一般向量丛上的则不能，这一限制正是附录 B 与附录 C 的起点。}

% CHUNK ch02e DONE
